\section{Which binary faces cause failure above the cutoff?}
A face fixing a positive fraction of coordinates governs the first-moment
cutoff, but the residual failures above it come from faces fixing at most two
coordinates.

\begin{corollary}\label{cor:rare-faces}
Take the binary model and error budget of Theorem~\ref{thm:box}, strictly
above the cutoff, and write \(K=e^{A n^{1/3}}\). Let \(F\) be the event that
some pair of faces is confusable, and let \(Z\) count all unordered confusable
pairs. Then
\[
 \prob(F)\sim\frac{n(n-1)(4E+1)}{\sqrt{\pi K}},
 \qquad \ex Z\sim 3\prob(F).
\]
For each coordinate \(i\), let \(U_i=\sum_{v\in S}(-1)^{v_i}\).
If we discard a landmark sample whenever two entries of the list
\((U_1,-U_1,\ldots,U_n,-U_n)\) are at distance at most \(4E\),
the discard probability is \((1+o(1))\prob(F)\). Conditional on retaining
the sample,
\[
 \lim_{n\to\infty}n^{-1/3}\log\prob(F\mid\text{retained})=-D,
 \qquad D=A-2B.
\]
\end{corollary}

\begin{proof}
We show that strictly above the cutoff, failure is dominated by pairs of faces
whose distance histograms differ along a single nonloop edge, and that
filtering out these one-edge collisions drops the failure rate to the two-edge
scale. We first rule out pairs of unequal, large, or multi-edge codimension
by union bounds, then classify all one-edge pairs and compute their
probabilities and pairwise overlaps by Fourier analysis, and finally prove the
matching conditional lower bound using a fixed codimension-three opposite pair.

\step{Eliminating high codimensions and multi-edge pairs.}
Put \(s=n^{1/3}\) and \(x=(E+1)/\sqrt K=e^{-Ds/2+o(s)}\).
For pairs of faces with different codimensions, the empirical-mean bound of
Section~\ref{sec:sum-codim} makes the summed failure probability
superexponentially small on the \(s\) scale.
For equal codimension \(r\ge\varepsilon n\), the strict first-moment margin
above the cutoff gives an exponentially small summed probability on the \(n\)
scale.
For intermediate equal codimensions \(1\le r<\varepsilon n\), the robust forest
estimate of Section~\ref{sec:robust-forest} bounds the confusability
probability of a single pair by
\[
 \exp[-c\Lambda\min(r,\sqrt{r\Lambda})],\qquad \Lambda=Ds .
\]
Take \(\varepsilon>0\) small. For \(\Lambda<r<\varepsilon n\), the logarithm
of the number of pairs is at most \(2r\log(2en/r)\), which is at most half the
exponent in absolute value once \(\varepsilon\) is small enough; summing over
this range costs at most \(\exp[-\Omega(n^{2/3})]\). For \(r\le \Lambda\), the
same pair count is absorbed by half of the exponent \(c\Lambda r\). Therefore,
for any fixed \(M\), choosing a sufficiently large fixed \(R\) makes the total
contribution from all codimensions above \(R\) equal to \(o(e^{-Ms})\).

Now consider pairs of equal codimension \(r\le R\). Draw an edge between two
distance levels for each distinct attainable unordered nonloop pair of
distances from a landmark to the two faces. Any two distinct edges form a
forest. Each chosen oriented edge is realized by some assignment of bits to the
union of the two face supports, defining a cell of density at least
\(2^{-2R}\). Condition on all landmarks outside these disjoint cells. Whenever
the two faces are confusable, each of the two forest cuts has an imbalance of
absolute value at most \(2E\), because the total positive histogram discrepancy
is at most that value. The two-cut bound therefore gives probability
\(O_R(x^2)\) for each pair with at least two edges. Summing over all such pairs
with \(r\le R\) gives \(O_R(n^{2R}x^2)\), which is negligible against \(n^2x\).

\step{Classifying one-edge pairs.}
It remains to classify pairs of equal codimension that have only one nonloop
edge. A landmark lying in the first face but not in the second produces some
edge \(\{0,k\}\); complementing all bits of that landmark produces the edge
\(\{r,r-k\}\), so having only one edge forces \(k=r\). Let \(b\) be the number
of conflicts between the two faces on their common support, and let \(c\) be
the size of each face's exclusive support. As a landmark varies over the first
face, its distance to the second face ranges over all integers from \(b\) to
\(b+c\). If \(c>0\) and \(r\ge2\), this range includes an interior distance
level, which produces a second edge and contradicts the one-edge assumption.
If \(c=0\), the two faces share the same support and must prescribe opposite
bits on every coordinate of that support. Their attainable nonloop pairs are
then \(\{j,r-j\}\), so there are \(\lceil r/2\rceil\) such edges. Requiring
only one edge thus implies that the one-edge pairs are precisely the pairs of
distinct codimension-one faces and the opposite pairs of codimension two with
common support.

For two codimension-one faces on different coordinates, the imbalance in the
single nonloop edge is, up to sign, \((U_i\pm U_j)/2\). An opposite pair on
the two-coordinate support \(\{i,j\}\) has the exact same imbalance, by the
identity
\(\mathbf1_{\text{both match}}-\mathbf1_{\text{both mismatch}}
=(\varepsilon_i+\varepsilon_j)/2\),
where \(\varepsilon_i,\varepsilon_j\) are the prescribed signs multiplied by
the respective coordinate signs. Thus, for every \(i<j\), each of the two
events
\[
 \big| (U_i\pm U_j)/2\big|\le 2E
\]
corresponds simultaneously to three confusable pairs: two pairs of
single-coordinate faces and one opposite pair of codimension two. Separately,
single-coordinate opposites on the same coordinate give the \(n\) events
\(|U_i|\le2E\), which contribute \(O(nx)\) in total and are negligible
compared to \(n^2x\).

\step{Probabilities and pairwise overlaps.}
For any dominant event, the linear form is the difference of two independent
\(\Bin(2^{n-2},K2^{-n})\) random variables, which has variance \(K(1-p)/2\)
with \(p=K2^{-n}\). On the fundamental frequency interval, its characteristic
function is \(\big(1-2p(1-p)(1-\cos t)\big)^{2^{n-2}}\). After rescaling by
the standard deviation, this characteristic function tends to the standard
Gaussian characteristic function and is dominated by \(e^{-c't^2}\) on the
rescaled interval. Furthermore, the Fourier phase for atom \(k\) tends to one
uniformly for \(|k|\le2E=o(\sqrt K)\) on compact frequency sets. By Fourier
inversion and dominated convergence, each dominant event has probability
\((1+o(1))(4E+1)/\sqrt{\pi K}\).

To bound pairwise overlaps, note that any two distinct dominant forms are
linearly independent functions of at most four landmark bits. We can therefore
choose two bit assignments with linearly independent contribution vectors and
condition on all landmarks outside their two disjoint cells. For each
permissible output pair of the two forms, at most one tuple of cell counts
works. Because the binomial maximal atom on each cell is \(O(K^{-1/2})\), the
intersection of two distinct dominant events costs \(O(x^2)\). Bonferroni's
inequalities then give an aggregate overlap of \(O(n^4x^2)=o(n^2x)\), and the
asymptotic equivalents for \(\prob(F)\) and \(\ex Z\) follow.

\step{The filter and conditional failure rate.}
The filtering rule removes precisely the one-edge events, and it can be checked
by sorting the signed values and comparing adjacent differences. Its rejection
probability follows from the same union estimate as above. Conditional on
retaining the sample, no one-edge pair is confusable, so the bounds for pairs
with two or more edges give a conditional log rate at most \(-D\).

For the matching lower bound, fix a pair of opposite faces on a
three-coordinate support, and write \(W_0,W_1\) for the signed count
differences in the disjoint shell pairs \((0,3),(1,2)\). These are independent
symmetric binomial differences, so the same Fourier estimate shows that
\(|W_0|,|W_1|\le E\) has probability \(\Theta(x^2)\) and implies confusability.
Let \(f_0,f_1\) be the corresponding single-landmark contribution functions.
Then \(3f_0+f_1\) is the sum of the three prescribed coordinate signs, and
\(f_0-f_1\) is their product. By orthogonality of the cube characters, no
rejected form, which is a nonzero linear form supported on at most two
coordinates, belongs to their span. Hence, for every rejected form, we can
choose three bit assignments with linearly independent contribution vectors on
the union of the at most five involved coordinates. Conditioning outside these
three disjoint cells and applying three maximal atom bounds gives an
intersection probability of \(O(x^3)\) with this fixed collision. Summing over
all rejected forms yields a total overlap of \(O(n^2x^3)=o(x^2)\), proving the
matching conditional lower rate.
\end{proof}
